\section{Модель, наблюдаемые и численная процедура}
\subsection{Геометрия и физические параметры}
Длинная ось сфероида направлена вдоль $z$; поверхность задаётся уравнением
$(x^2+y^2)/a^2+z^2/c^2=1$. Отношение полуосей равно 2,7778.
Диэлектрическая проницаемость однородной среды $\varepsilon_m=2{,}25$.
Радиус КТ $r_q=2$ нм используется при вычислении расстояния между
центрами; оптический переход КТ представлен точечным диполем.
Зазор $g$ измеряется между поверхностями, поэтому расстояние от
поверхности металла до центра диполя равно $r_q+g$.
Рассмотрены $g=0{,}5;\ 0{,}75;\ 1;\ 1{,}5;\ 2;\ 3;\ 5;\ 7;\ 10;\ 12$ нм.

\begin{table}[!htbp]\centering\small
\caption{Геометрические каналы. $R$ --- расстояние между центрами;
$G$ --- угловой множитель дипольного приближения.}
\label{tab:geometry}
\begin{tabular}{p{4.0cm}lll}
\toprule
Канал & Положение КТ & Поляризация & $G$\\\midrule
Торцевой продольный & $(0,0,c+r_q+g)$ & $z$ & $2$\\
Торцевой поперечный & $(0,0,c+r_q+g)$ & $x$ & $-1$\\
Боковой продольный & $(a+r_q+g,0,0)$ & $z$ & $-1$\\
Боковой радиальный & $(a+r_q+g,0,0)$ & $x$ & $2$\\
Боковой тангенциальный & $(a+r_q+g,0,0)$ & $y$ & $-1$\\\bottomrule
\end{tabular}
\end{table}
Два поперечных боковых канала различаются направлением диполя относительно
радиуса-вектора и сохраняются раздельно во всех сравнениях.
При минимальном зазоре $R=15$ нм у торца и 7 нм сбоку.

Энергия перехода КТ $\hbar\omega_0=2{,}042$ эВ, эффективный дипольный
момент $d=13{,}9$ D, энергия распада заселённости
$\hbar\gamma_1=268$ нэВ и чистого дефазирования
$\hbar\gamma_\varphi=1{,}27$ мэВ. Эти параметры и длительность импульса
соответствуют набору, использованному для КТ в~\cite{shah2013}.
При этом здесь изучается \emph{одна} металлическая частица и не вводится
поправка к $\gamma_1$, имитирующая эффект Парселла в той работе.
Полная скорость затухания когерентности
$\Gamma_2=\gamma_1/2+\gamma_\varphi$ соответствует энергии
1,270134 мэВ: $T_1=2{,}456$ нс, $T_2=0{,}5182$ пс.
Используется соглашение об эффективном внешнем диполе КТ, с множителем
локального поля $f_{\rm loc}=1$. Дополнительное умножение на
$3\varepsilon_m/(\varepsilon_q+2\varepsilon_m)$ означало бы другое
соглашение и в данном расчёте не выполняется.

\subsection{Электростатический отклик металла и обратная связь}
Далее формулы для поляризуемостей и уравнения динамики записаны в атомных
единицах, если не указано иное. Для данного направления диполя отклик
наночастицы удобно представить тремя частотными функциями:
\begin{equation}
p_m=A E_{\rm inc}+B p_q,\qquad
E_m=B E_{\rm inc}+K p_q.\label{eq:ABK}
\end{equation}
Здесь $p_m$ --- полный диполь металла, $p_q$ --- диполь КТ,
$E_m$ --- поле металла в центре КТ. Равенство двух перекрёстных
коэффициентов $B$ выражает взаимность при согласованных проекциях.

В дипольном приближении (DD)
\begin{equation}
A=C H_{L_1},\quad B=AJ,\quad K=AJ^2,\qquad
C=\frac{\varepsilon_m a^2c}{3},\quad
J=\frac{G}{\varepsilon_m R^3},\label{eq:DD}
\end{equation}
\begin{equation}
H_L(\omega)=\frac{\varepsilon_{\rm Au}(\omega)-\varepsilon_m}
{\varepsilon_m+L[\varepsilon_{\rm Au}(\omega)-\varepsilon_m]}.
\label{eq:HL}
\end{equation}
Величины $A$ и $K$ имеют размерности объёма и обратного объёма,
а $B$ безразмерен. Для полного квазистатического расчёта (FQS)
используется решение уравнения Лапласа в сфероидальных гармониках.
В нём $A$ остаётся откликом яркой дипольной моды, $B$ определяется её
точным внешним полем, а реакционное поле включает высшие пространственные
моды:
\begin{equation}
K(\omega)=\sum_\nu w_\nu H_{L_\nu}(\omega),\qquad w_\nu\geq0.
\label{eq:kernel}
\end{equation}
Геометрические коэффициенты $L_\nu$ и $w_\nu$ вычисляются из
сфероидальных функций и положения источника. Поэтому различие DD/FQS
включает как поправку к возбуждающему полю $B$, так и изменение
самодействия $K$. FQS означает полноту пространственного ряда в
квазистатике, а не решение волновых уравнений Максвелла.

Дисперсия золота задаётся таблицей Джонсона--Кристи~\cite{johnson1972}:
интерполируются $n$ и $k$, затем вычисляется
$\varepsilon_{\rm Au}=(n+ik)^2$.
Для временной задачи яркий отклик аппроксимируется пассивной суммой
\begin{equation}
H_{L_1}(\omega)=\alpha_\infty+
\sum_{j=1}^{N}\frac{f_j}{\omega_j^2-\omega^2-i\gamma_j\omega},
\qquad f_j\geq0,\quad\gamma_j>0.\label{eq:lorentz}
\end{equation}
Константа $\alpha_\infty$ фиксируется условием
$\varepsilon_{\rm Au}(\omega\to\infty)=1$ и не подгоняется свободно.
Остальные моды получаются из \emph{того же} материала посредством
$H_L=H_{L_1}/[1+(L-L_1)H_{L_1}]$. Основная аппроксимация содержит
$N=12$ лоренцевых осцилляторов в окне $0{,}8$--$3{,}5$ эВ;
$N=13$ используется в контроле, $N=1$ --- как намеренно грубое сравнение.
Каждое лоренцево слагаемое имеет два комплексных частотных полюса.
Число осцилляторов $N$ не равно максимальной степени пространственного ряда $n_{\max}$.

\subsection{Импульсное возбуждение}
Динамика КТ описывается компонентами $W=\rho_{ee}-\rho_{gg}$,
$P=2\operatorname{Re}\rho_{ge}$, $Q=-2\operatorname{Im}\rho_{ge}$:
\begin{align}
\dot W&=\Omega Q-\gamma_1(W+1),\nonumber\\
\dot Q&=-\omega_0P-\Omega W-\Gamma_2Q,\nonumber\\
\dot P&=\omega_0Q-\Gamma_2P,\qquad
\Omega=2d(E_{\rm inc}+E_m),\quad p_q=dP.
\label{eq:bloch}
\end{align}
В начальный момент $W=-1$, $P=Q=0$, металл невозбуждён.
Уравнения интегрируются с действительным полем, без приближения вращающейся
волны. Лоренцевы слагаемые реализуются затухающими вспомогательными
осцилляторами. Для тёмной части пространственного спектра используется
положительная редукция с отдельными проверками точности.

Падающий импульс имеет вид
\begin{equation}
E_{\rm inc}(t)=E_0e^{-t^2/(2\sigma^2)}\cos\omega_Lt,
\qquad \hbar\omega_L=2{,}042\ \text{эВ}.
\end{equation}
Полная ширина \emph{интенсивности} на половине максимума равна 20 фс,
поэтому $\sigma=12{,}011$ фс. Спектральная ширина интенсивности составляет
0,09125 эВ. В СИ флюенс определяется как
\begin{equation}
\mathcal F=I_0\sqrt\pi\sigma,\qquad
I_0=\tfrac12\sqrt{\varepsilon_m}\varepsilon_0c_0E_0^2,
\label{eq:fluence}
\end{equation}
где $c_0$ --- скорость света. Эта нормировка учитывает среду.
Заселённость считывается при $t_{\rm read}=3{,}6$ пс относительно центра
импульса. Собственный распад за этот интервал составляет 0,1465\%,
но когерентный отклик и поле металла успевают затухнуть в пределах
заданных критериев.

\subsection{Различаемые наблюдаемые}
Линейная поляризуемость изолированной КТ и её гибридный отклик равны
\begin{equation}
\beta(\omega)=\frac{2d^2\omega_0}{\omega_0^2+(\Gamma_2-i\omega)^2},
\qquad \frac{p_q}{E_{\rm inc}}=
\frac{\beta(1+B)}{1-\beta K},\qquad
S(E)=\left|\frac{p_q}{E_{\rm inc}}\right|^2.\label{eq:linear}
\end{equation}
$S$ --- квадрат дипольного отклика, а не заселённость или интенсивность
фотолюминесценции. Используются два усиления:
\begin{equation}
G_{\rm fixed}=\frac{S(E_q)}{S_0(E_q)},\qquad
G_{\rm opt}=\frac{\max_{E\in\mathcal W}S(E)}{\max_{E\in\mathcal W}S_0(E)},
\quad E_q=\hbar\omega_0=2{,}042\ \text{эВ},
\end{equation}
где $S_0$ относится к изолированной КТ,
$\mathcal W=[1{,}872;2{,}212]$ эВ. Основная спектральная сетка содержит
4001 точку от 1,85 до 2,234 эВ. Ширина $\Gamma_0=2{,}54142$ мэВ
изолированной линии определяется тем же алгоритмом, что и гибридная:
по половине выраженности пика над его локальным основанием.
При конкурирующих пиках единственная ширина не назначается.
Ширина плазмонной полосы материала, напротив, извлекается относительно
глобального основания; эти определения не взаимозаменяемы.

Спектральная ошибка DD относительно FQS вычисляется без нормирования
каждой кривой на собственный максимум:
\begin{equation}
\delta_{\rm spec}=\left[
\frac{\int_{\mathcal W}(S_{\rm DD}-S_{\rm FQS})^2dE}
{\int_{\mathcal W}S_{\rm FQS}^2dE}\right]^{1/2}.
\end{equation}
Нелинейный порог $\Fth$ соответствует первому достижению
$\rho_{ee}(t_{\rm read})=0{,}5$ на первой возрастающей ветви зависимости
от $\sqrt{\mathcal F}$. Уточняется заключённый в скобки корень;
позднейшие пересечения после рабиевского максимума не подменяют первый
порог. Ошибка порога
$\delta_F=|\Fth^{\rm DD}/\Fth^{\rm FQS}-1|$ определяется лишь при
разрешённых порогах обеих моделей. Отсутствие порога до верхнего предела
$200\ \uJ$ фиксируется отдельно от случая, когда первая рабиевская
ветвь вообще не достигает уровня 0,5.

Для дополнительного спектра S1 преобразуется полный диполь $p_q+p_m$:
\begin{equation}
\alpha_{\rm eff}(\omega;\mathcal F)=
\frac{\widetilde p_q+\widetilde p_m}{\varepsilon_m\widetilde E_{\rm inc}},
\qquad
\sigma_{\rm QS,work}=\frac{k}{\varepsilon_0}
\operatorname{Im}\alpha_{\rm eff}^{\rm SI},\quad
k=\sqrt{\varepsilon_m}\omega/c_0.\label{eq:work}
\end{equation}
Здесь используется преобразование $\widetilde f=\int f(t)e^{i\omega t}dt$;
поляризуемость в последнем равенстве переводится в СИ.
Это квазистатический показатель спектральной работы поля. Для сильного
импульса он зависит от флюенса и не является линейной восприимчивостью.
Его не отождествляют с чистым поглощением или рассеянием.
Разность $\Delta\sigma=\sigma_{\rm hybrid}-\sigma_{\rm bare\,metal}$
берётся внутри одной материальной ветви. Знак этой разности не задаёт
знак полной работы $W_{\rm ext}=\int E_{\rm inc}\dot p_{\rm total}dt$.

\subsection{Контроль численного решения}
Основные расчёты выполнены при $n_{\max}=80$ и методом DOP853 с
$\mathrm{rtol}=10^{-7}$, $\mathrm{atol}=10^{-9}$; учитывается ограничение
шага по возбуждаемой частотной полосе. Контролируются пространственная
сходимость, модальная аппроксимация, устойчивость связанной линейной
системы, положительность матрицы плотности, спектральная утечка и
затухание временного хвоста. Для $N=12$ предельные ошибки яркого отклика
составляют 2,5\% по нормированному RMS и 5\% по максимальной
относительной ошибке. Для пространственных мод допуски номинального
сценария равны 10 и 20\%, а конечные наблюдаемые дополнительно сверяются
с прямым табличным материалом. Допуск NRMS спектра КТ равен 5\%,
сдвига --- $0{,}1\Gamma_0$, ширины и усиления --- 10\%.
Здесь нормированная среднеквадратичная ошибка (NRMS) комплексного
отклика $X$ на соответствующей контрольной сетке $E_i$ означает
\begin{equation}
\mathrm{NRMS}(X)=
\frac{[\sum_i|X_{\rm fit}(E_i)-X_{\rm ref}(E_i)|^2]^{1/2}}
{[\sum_i|X_{\rm ref}(E_i)|^2]^{1/2}}.
\end{equation}
Для материальной подгонки отдельно проверяются $X=\alpha$ и
$X=1/\alpha$, для конечного спектра --- $X=S$.
Максимальная относительная материальная ошибка равна
$\max_i|\alpha_{\rm fit}(E_i)-\alpha_{\rm ref}(E_i)|/
|\alpha_{\rm ref}(E_i)|$.
Результаты фактических проверок приведены в разделе~\ref{sec:validation}.
